% Part: computability
% Chapter: recursive-functions
% Section: general-recursive-functions

\documentclass[../../../include/open-logic-section]{subfiles}

\begin{document}

\olfileid{cmp}{rec}{gen}
\olsection{عمومي بازګشتي تابعې}

د هرځاے تعريف شويو تابعو د يو سټ د تر لاسه کولو بله لار هم شته۔
هرځاے تعريف شوې تابع $f(x,\vec z)$ هغه وخت \emph{منظمه} وبولئ چې د
طبيعي عددونو د هرې لړۍ~$\vec z$ دپاره يو~$x$ موجود وي چې
$f(x,\vec z)=0$ وي۔ په بله وينا، منظمې تابعې کټ مټ هغه تابعې دي چې
بې‌حده پلټنه پرې عملي کېدے شي او په پايله کښې يوه هرځاے تعريف شوې
تابع تر لاسه کېږي۔ بې‌حده پلټنه په محتاط ډول يوازې منظمو تابعو ته
محدودولے شو:

\begin{defn}
\ollabel{defn:general-recursive}
د \emph{عمومي بازګشتي تابعو} سټ له طبيعي عددونو څخه طبيعي عددونو ته د
بېلابېلو ځاييزو کچو د تابعو تر ټولو کوچنے سټ دے چې صفر، جانشين او
پروجکشنونه لري، او د ترکيب، بنسټيز بازګښت او پر \emph{منظمو} تابعو د
عملي شوې بې‌حده پلټنې لاندې بند دے۔
\end{defn}

په څرګند ډول هره عمومي بازګشتي تابع هرځاے تعريف شوې ده۔ د
\olref{defn:general-recursive} او \olref[par]{defn:recursive-fn} تر
منځ توپير دا دے چې په وروستي تعريف کښې د جوړونې په منځ کښې د جزوي
بازګشتي تابعو کارول هم روا دي؛ يوازينے شرط دا دے چې په پاے کښې تر لاسه
شوې تابع هرځاے تعريف شوې وي۔ نو د ``عمومي'' ټکی، چې يو تاريخي پاتې
شونی دے، سم نوم نۀ دے؛ په ظاهره \olref{defn:general-recursive} تر
\olref[par]{defn:recursive-fn} \emph{لږ} عمومي دے۔ خو له نېکه مرغه دا
توپير ظاهري دے؛ که څه هم تعريفونه بېل دي، د عمومي بازګشتي تابعو سټ او
د بازګشتي تابعو سټ يو او هماغه سټ دے۔

\end{document}
